\begin{apendicesenv}
\partapendices

\chapter{Roteiro e Síntese da Entrevista com Lucas Mathias}
\label{apend:entrevista}

Este apêndice apresenta o roteiro de perguntas utilizado na entrevista estruturada com Lucas Mathias, PM de experimentos da Capim Software, realizada por videoconferência em 7 de maio de 2026, organizado por tema junto a uma síntese das respostas obtidas. A transcrição integral e literal da conversa está disponível no repositório do trabalho, no arquivo \texttt{report/transcript\_enterview.md}.

\section{Ferramentas de trabalho}
\textbf{Pergunta:} Quais softwares a empresa utiliza para gerenciar seus projetos?\\
\textbf{Síntese:} Notion, Cursor/Claude (desenvolvimento assistido por IA), Linear (gestão de tarefas) e N8N (automação). Mixpanel e Metabase são usados para acompanhamento de métricas, mas não para construção de produto.

\section{Trajetória competitiva}
\textbf{Pergunta:} Qual o principal diferencial da Capim frente aos concorrentes? Como a empresa se inseriu no mercado?\\
\textbf{Síntese:} A Capim não precisou se diferenciar ativamente no início, pois criou uma categoria que não existia (crédito odontológico no ponto de venda). O diferencial evoluiu produto a produto: pioneirismo no crédito, simplicidade e modernidade no sistema de gestão, e competição por taxa na maquininha.

\section{Serviços sob constante aprimoramento}
\textbf{Pergunta:} Quais serviços ou ferramentas estão sob constante revisão?\\
\textbf{Síntese:} Praticamente todos os produtos, inclusive o sistema de gestão -- considerado internamente menos prioritário, mas ainda alvo de melhorias incrementais (ex.: lembrete automático de aniversário de pacientes). Exemplos de inovação incremental: Capix (fluxo de crédito via Pix, antes inteiramente operado por WhatsApp). Exemplo de inovação radical: CamilA, assistente de IA para vender crédito diretamente ao paciente, com a Capim atuando como intermediária entre paciente e clínica (modelo descrito pelo entrevistado como ``Uber dos dentistas'').

\section{Exploração de novos mercados}
\textbf{Pergunta:} A empresa planeja explorar mercados além da odontologia?\\
\textbf{Síntese:} A Capim não começou como uma empresa de odontologia; chegou a testar diversificação para outras áreas de saúde (ótica, clínicas médicas, cirurgias eletivas) antes de se especializar em odontologia por absorção de demanda. A liderança considera que ainda há espaço relevante para crescer dentro do nicho atual antes de voltar a diversificar.

\section{Verba e equipe dedicadas à inovação}
\textbf{Pergunta:} Há verba e equipe específicas para projetos de inovação?\\
\textbf{Síntese:} Cada squad é responsável por seu produto, com uma coordenação central para decisões de portfólio. A maior parte das equipes -- inclusive o squad de sustentação -- realiza inovações incrementais, ainda que com intensidade variável.

\section{Pipeline de aprovação e avaliação}
\textbf{Pergunta:} Há um pipeline formal para aprovar e avaliar novos projetos?\\
\textbf{Síntese:} O desenvolvimento envolve progressivamente mais pessoas e equipes, controlado pelo número de clínicas que conseguem visualizar a solução (de aproximadamente 10, para 100, até cerca de 500 clínicas). O critério de sucesso de cada iniciativa é definido no início e pode ser revisado ao longo do projeto -- como ocorreu no projeto B2C, cuja métrica mudou de ``contratos assinados'' para ``tratamentos fechados pelas clínicas''. Projetos sem sucesso são descontinuados e documentados, ou ``congelados'' aguardando maturação tecnológica -- caso da CamilA.

\section{Gestão de incertezas}
\textbf{Pergunta:} Como é feita a gestão de incertezas?\\
\textbf{Síntese:} Cada projeto recebe uma métrica de sucesso própria; o time persiste na iniciativa até o fim do ciclo trimestral antes de decidir por descontinuação. O projeto da CamilA foi mantido por três meses antes da decisão de congelamento.

\section{Avaliação de projetos com potencial latente}
\textbf{Pergunta:} Como são avaliados projetos com potencial ainda não comprovado?\\
\textbf{Síntese:} A maior oportunidade percebida atualmente é a aquisição de pacientes pelas clínicas (``agenda cheia e rentável''), validada por meio de conversas diretas com clínicas parceiras sobre dores reais de ocupação de agenda.

\section{Oportunidades e ameaças}
\textbf{Pergunta:} Quais são as maiores oportunidades e ameaças da Capim?\\
\textbf{Síntese:} A principal ameaça percebida é a dependência do cenário macroeconômico -- em particular o nível de endividamento da população, que limita a aprovação de crédito. O entrevistado resume o princípio de inovação da empresa como ``iterar o mais rapidamente possível, colocar na mão do usuário para ele testar o mais rápido possível e modificar o produto o mais rápido possível''.

\end{apendicesenv}
